\documentclass[11pt,a4paper]{article}
\usepackage[margin=23mm]{geometry}
\usepackage{amsmath,amssymb,bm}
\usepackage[T1]{fontenc}
\usepackage[hypertexnames=false]{hyperref}
\newcommand{\wrap}{\operatorname{wrap}}
\newcommand{\diag}{\operatorname{diag}}
\newcommand{\clip}{\operatorname{clip}}
\title{HOFA-MPC: Equations, Derivations and Identification}
\author{Implementation notes for the current workspace}
\date{12 October 2026}
\begin{document}
\maketitle
These equations document the current three-DOF controller. The accompanying
Chinese Markdown file explains the source evidence and historical validation.
Equations E16 and E17 deliberately distinguish the consistent correction-input
model from the current implementation. No controller changes are implied.

\section{Coordinates and dynamics}
Let $\eta=[x,y,\psi]^T$, $\nu=[u,v,r]^T$, $\tau=[X,Y,N_z]^T$,
and $\xi=[\eta^T,\nu^T]^T$. The controller uses ENU/FLU internally.
\begin{align}
x_E&=y_N,&y_E&=x_N,&\psi_E&=\wrap(\pi/2-\psi_N),\notag\\
u_{\mathrm{FLU}}&=u_{\mathrm{FRD}},&v_{\mathrm{FLU}}&=-v_{\mathrm{FRD}},&
r_{\mathrm{FLU}}&=-r_{\mathrm{FRD}},\notag\\
\tau_{\mathrm{FRD}}&=\diag(1,-1,-1)\tau_{\mathrm{FLU}}.\tag{E01}
\end{align}
\begin{equation}
\dot\eta=J(\psi)\nu,\qquad
M\dot\nu+C(\nu)\nu+D(\nu)\nu=\tau+d_\tau.\tag{E02}
\end{equation}
\begin{equation}
J=\begin{bmatrix}\cos\psi&-\sin\psi&0\\\sin\psi&\cos\psi&0\\0&0&1\end{bmatrix},\qquad
\dot J=r\begin{bmatrix}-\sin\psi&-\cos\psi&0\\\cos\psi&-\sin\psi&0\\0&0&0\end{bmatrix}.
\tag{E03}
\end{equation}
\begin{align}
M&=M_{\mathrm{dry}}+M_A=\diag(m_x,m_y,I_z),\notag\\
D(\nu)&=\diag(d_{l,x}+d_{q,x}|u|,d_{l,y}+d_{q,y}|v|,d_{l,\psi}+d_{q,\psi}|r|).
\tag{E04}
\end{align}
For diagonal inertia, the implemented optional Coriolis matrix is
\begin{equation}
C(\nu)=\begin{bmatrix}0&0&-m_yv\\0&0&m_xu\\m_yv&-m_xu&0\end{bmatrix}.
\tag{E05}
\end{equation}
The current simulation configuration disables this matrix, i.e., $C=0$.
The kinematic term $\dot J\nu$ is still retained.

\section{HOFA transformation}
Differentiating the kinematics and substituting the nominal dynamics gives
\begin{align}
\ddot\eta&=\dot J\nu+J\dot\nu
=\dot J\nu-JM^{-1}(C\nu+D\nu)+JM^{-1}\tau\notag\\
&=b(\xi)+G_\tau(\xi)\tau,\notag\\
b(\xi)&=\dot J\nu-JM^{-1}(C\nu+D\nu),\qquad G_\tau=JM^{-1}.
\tag{E06}
\end{align}
Since $J^{-1}=J^T$, imposing $\ddot\eta=a_c$ yields
\begin{equation}
\boxed{\tau_d=MJ^T(a_c-\dot J\nu)+C\nu+D\nu.}\tag{E07}
\end{equation}
This is exact for the nominal model with exact force realization. With an
additive force disturbance and otherwise exact modeling,
$\ddot\eta=a_c+JM^{-1}d_\tau$.

\section{Time-indexed reference}
Let $s$ denote arc length and $V=\dot s$ denote path speed (distinct from sway
velocity $v$). Then
\begin{equation}
\kappa=\frac{d\psi_r}{ds},\qquad \dot\psi_r=V\kappa,\qquad
\ddot\psi_r=\dot V\kappa+V^2\kappa',\qquad \kappa'=\frac{d\kappa}{ds}.
\tag{E08}
\end{equation}
The implemented speed heuristic is
\begin{align}
a_{\ell,\max}&=\max(a_{\mathrm{acc}},a_{\mathrm{dec}}),\notag\\
b_\psi&=\max\{\alpha_{\max}-|\kappa|a_{\ell,\max},\rho\alpha_{\max}\},\notag\\
V_{\lim}&=\min\left\{V_{\max},\frac{r_{\max}}{\max(|\kappa|,\epsilon)},
\sqrt{\frac{b_\psi}{\max(|\kappa'|,\epsilon)}}\right\}.\tag{E09}
\end{align}
The budget floor prevents artificial stops; it relaxes the worst-case yaw
acceleration argument and is not a universal hard-bound proof.
Forward/backward passes and the approximate arrival-time table are
\begin{align}
V_i&\leftarrow\min(V_i,\sqrt{V_{i-1}^2+2a_{\mathrm{acc}}\Delta s_i}),\notag\\
V_i&\leftarrow\min(V_i,\sqrt{V_{i+1}^2+2a_{\mathrm{dec}}\Delta s_{i+1}}),
\qquad V_{\mathrm{end}}=0,\notag\\
T_i&=T_{i-1}+\frac{\Delta s_i}{\max(0.08,(V_{i-1}+V_i)/2)}.\tag{E10}
\end{align}
Current settings are $V_{\max}=0.20$, $r_{\max}=0.5$, $\alpha_{\max}=0.25$,
$a_{\mathrm{acc}}=0.12$, $a_{\mathrm{dec}}=0.18$, $\rho=0.2$ in SI units.

\section{Layer 1: feasible acceleration boxes}
For horizontal thruster $j$, its position and direction determine
\begin{equation}
B_{h,:,j}=\begin{bmatrix}d_{x,j}\\d_{y,j}\\p_{x,j}d_{y,j}-p_{y,j}d_{x,j}\end{bmatrix},\qquad
\tau=B_hf,\qquad \underline f\le f\le\overline f.\tag{E11}
\end{equation}
At a frozen nominal state,
\begin{align}
a&=b+Hf,\qquad H=JM^{-1}B_h,\notag\\
\mathcal A(\xi)&=\{b+Hf:\underline f\le f\le\overline f\},\notag\\
f_c&=(\overline f+\underline f)/2,\qquad f_\Delta=(\overline f-\underline f)/2,\notag\\
c&=b+Hf_c,\qquad\bar h_i=\sum_j|H_{ij}|f_{\Delta,j}.\tag{E12}
\end{align}
The outer projection box need not be jointly reachable. The implemented LP
constructs an inner box using one feasible force witness per corner:
\begin{equation}
\begin{aligned}
\max_{h,\beta,\{f^{(s)}\}}\quad &\beta+10^{-4}\sum_{i\in\mathcal I}h_i/\bar h_i\\
\text{subject to}\quad
&Hf^{(s)}=Hf_c+\diag(s)h,\quad s\in\{-1,1\}^3,\\
&\underline f\le f^{(s)}\le\overline f,\\
&0\le h_i\le\gamma\bar h_i,\quad0\le\beta\le\gamma,\\
&\beta\bar h_i\le h_i\quad(i\in\mathcal I),\quad
\mathcal I=\{i:\bar h_i>10^{-12}\},\\
&\underline a=c-h,\qquad\overline a=c+h.
\end{aligned}\tag{E13}
\end{equation}
Here $\gamma=1/3$. There are 36 variables for four horizontal thrusters.
Convexity of the reachable set implies containment of the entire box once all
eight corners have feasible witnesses. This does not establish robustness to
state prediction errors or exact realization by the downstream pseudoinverse.

\section{Layer 2: error MPC}
\begin{equation}
e=\eta-\eta_r,\quad e_\psi=\wrap(\psi-\psi_r),\quad
\dot e=J\nu-\dot\eta_r,\quad z=\begin{bmatrix}e\\\dot e\end{bmatrix},\quad
a_c=a_r+w.\tag{E14}
\end{equation}
For zero-order-held acceleration with $h=0.1$ s,
\begin{equation}
A_d=\begin{bmatrix}I_3&hI_3\\0&I_3\end{bmatrix},\qquad
B_d=\begin{bmatrix}\frac12h^2I_3\\hI_3\end{bmatrix}.\tag{E15}
\end{equation}
Let $q_r=[\eta_r^T,\dot\eta_r^T]^T$. Subtracting the reference from the
physical double integrator gives the consistent correction-input model
\begin{align}
z_{i+1}&=A_dz_i+B_dw_i+\delta_i^{\mathrm{consistent}},\notag\\
\delta_i^{\mathrm{consistent}}&=A_dq_{r,i}+B_da_{r,i}-q_{r,i+1}.\tag{E16}
\end{align}
The current source instead implements
\begin{align}
\delta_i^{\mathrm{code}}&=A_dq_{r,i}-q_{r,i+1},\quad0\le i<N_p-1,\notag\\
\delta_{N_p-1}^{\mathrm{code}}&=0,\notag\\
z_{i+1}&=A_dz_i+B_dw_i+\delta_i^{\mathrm{code}},\notag\\
\delta_{i,\psi}^{\mathrm{code}}&=\wrap(h\dot\psi_{r,i}+\psi_{r,i}-\psi_{r,i+1}).
\tag{E17}
\end{align}
Yet it shifts bounds and adds output feedforward according to the correction
interpretation:
\begin{equation}
\underline a_i-a_{r,i}\le w_i\le\overline a_i-a_{r,i},\qquad
a_c=a_{r,0}+w_0^*.\tag{E18}
\end{equation}
Consequently E16 and E17 differ for nonzero reference acceleration. For an
exactly discretized accelerated reference, $z_0=w_0=0$ gives zero next-step
error under E16 but $-B_da_{r,0}$ under E17. This document records the mismatch
without silently changing the implementation.

The implemented objective, with $N_p=20$, is
\begin{align}
\min_{w_0,\ldots,w_{N_p-1}}\quad\mathcal J
&=\sum_{i=0}^{N_p-1}\left(z_i^TQ_i z_i+w_i^TRw_i+\Delta w_i^TS\Delta w_i\right),\notag\\
Q_i&=\begin{cases}Q&i<N_p-1,\\F&i=N_p-1,\end{cases}\quad F=4Q,\notag\\
\Delta w_0&=w_0,\qquad\Delta w_i=w_i-w_{i-1}\quad(i\ge1).\tag{E19}
\end{align}
The constraints are E17--E18. Current weights are
\begin{align*}
Q&=\diag(20,20,20,10,10,8),\\
R&=\diag(0.20,0.20,0.14),\qquad S=\diag(1.20,1.20,0.75).
\end{align*}
The terminal weight acts on $z_{N_p-1}$, not $z_{N_p}$. The first increment
is relative to zero, not to the previously applied command. With frozen bounds
this is a convex quadratic box-constrained problem, solved by L-BFGS-B.

\section{Adjoint derivative and computational changes}
Consider
\[
\mathcal L=\mathcal J+\sum_{i=0}^{N_p-1}\lambda_{i+1}^T
(A_dz_i+B_dw_i+\delta_i-z_{i+1}).
\]
Differentiating with respect to intermediate states and controls gives
\begin{align}
\lambda_{N_p}&=0,\notag\\
\lambda_i&=2Q_i z_i+A_d^T\lambda_{i+1},\notag\\
\nabla_{w_i}\mathcal J&=2Rw_i+B_d^T\lambda_{i+1}+2S\Delta w_i
-\mathbf1_{i<N_p-1}\,2S\Delta w_{i+1}.\tag{E20}
\end{align}
One forward rollout and one backward sweep cost $O(N_p)$ at fixed state
dimension. Finite differences over $3N_p$ variables cost $O(N_p^2)$ per full
gradient. This statement concerns evaluations, not the total optimization
iteration complexity. The shifted warm start is
\begin{equation}
W_k^{(0)}=[(w_{k-1,1}^*)^T,\ldots,(w_{k-1,N_p-1}^*)^T,0^T]^T.\tag{E21}
\end{equation}
The current objective returns cost and gradient together with
\texttt{jac=True}, avoiding duplicate evaluations at the same iterate.
The saved offline solve-only benchmark reports:
\begin{center}
\begin{tabular}{lrr}
\hline
Latency & Before joint return & After joint return\\\hline
Mean (ms)&26.297&13.934\\
Median (ms)&25.060&13.187\\
P95 (ms)&36.701&19.382\\\hline
\end{tabular}
\end{center}
The mean speedup is $1.887\times$ (47.013\% reduction). Both versions already
used analytic gradients. The benchmark used 120 reconstructed cases, three
timed rounds, and 360 solves per version. Paired full sequences and objectives
had zero maximum difference. This excludes Layer 1 and ROS callback costs;
it is not an online end-to-end speedup claim.

\section{Allocation and calibrated force commands}
The library-level four-thruster allocator computes
\begin{equation}
f_{\mathrm{lib}}=\clip\left((B_h^TW_aB_h+\lambda_a I)^{-1}B_h^TW_a\tau,
\underline f,\overline f\right),\quad
W_a=\diag(1,1,2),\quad\lambda_a=10^{-4}.\tag{E22}
\end{equation}
The running launch path instead uses an external six-thruster allocator,
including depth actuation, in the FRD frame:
\begin{equation}
\tau_4=[X,Y,N_z,Z]^T,\qquad
f_6=A_6^T(A_6A_6^T)^{-1}\tau_4.\tag{E23}
\end{equation}
Positive and negative thrust curves are inverted separately by piecewise
linear interpolation:
\begin{equation}
c(F)=c_l+\frac{|F|-F_l}{F_u-F_l}(c_u-c_l),\qquad
\mathrm{PWM}=\operatorname{sign}(F)c(F).\tag{E24}
\end{equation}
The code also zeros forces of magnitude at most 0.05 N and saturates at curve
endpoints. Pseudoinverse allocation followed by saturation is not an exact
bounded least-squares optimizer.

\section{Two-stage identification and validation}
For one axis, use $F=m_{\mathrm{eff}}\dot\chi+d_l\chi+d_q|\chi|\chi$,
where $\chi$ is $u$, $v$, or $r$. In yaw, inertia has units kg\,m$^2$.
Steady sweeps identify drag first; fixed-drag transients then identify inertia:
\begin{align}
\Phi_{j,:}&=[\bar\chi_j,|\bar\chi_j|\bar\chi_j],\qquad
y_j=\overline{F_j-m_{\mathrm{given}}\dot\chi_j},\notag\\
[\hat d_l,\hat d_q]^T&=\arg\min_{\theta\ge0}\|\Phi\theta-y\|_2^2,\notag\\
b_k&=F_k-\hat d_l\chi_k-\hat d_q|\chi_k|\chi_k,\notag\\
\hat m_{\mathrm{eff}}^{\mathrm{derivative}}
&=\frac{\sum_k\dot\chi_kb_k}{\sum_k\dot\chi_k^2},\qquad
\hat m_A=\hat m_{\mathrm{eff}}-m_{\mathrm{dry}},\notag\\
\dot{\hat\chi}(t;m)&=\frac{F(t)-\hat d_l\hat\chi-\hat d_q|\hat\chi|\hat\chi}{m},\notag\\
\hat m_{\mathrm{eff}}^{\mathrm{response}}
&=\arg\min_{m>0}\sum_{j,k}[\hat\chi_j(t_k;m)-\chi_j(t_k)]^2.\tag{E25}
\end{align}
The drag estimate is mass-independent only in the steady-state limit.
Different nonzero speed magnitudes are required to separate the two drag
terms; sign reversal alone at one magnitude is insufficient.
The derivative estimator follows from differentiating a one-parameter
least-squares objective. The response estimator integrates the scalar ODE,
avoiding numerical differentiation of measured speed. Historical reports
document response fitting; the current added-mass CLI implements the
derivative estimator.

The current Stonefish parameters are
\begin{align*}
M_{\mathrm{dry}}&=\diag(7.95,7.95,0.174887),\\
M_A&=\diag(26.254017,26.254017,0.067653),\\
M&=\diag(34.204017,34.204017,0.242540),\\
d_l&=[35.488892,35.656137,1.163915]^T,\\
d_q&=[17.853671,17.584671,0.021297]^T.
\end{align*}
Each axis used ten signed steady trials. Surge used force magnitudes
4, 6, 8, 10, 12 N for 25 s each; sway used 3, 4, 5, 6, 8 N for 20 s;
yaw used 0.2, 0.3, 0.4, 0.5, 0.6 N\,m for 20 s.

The accepted translational inertia is the current simulator geometry value.
On the selected validation transients, surge RMSE was 0.002222 m/s using
34.204 kg versus 0.005023 m/s using the step-fitted 41.184 kg. Sway RMSE was
0.001571 m/s using 34.204 kg versus 0.004800 m/s using 43.833 kg.
The validation subsets excluded both lowest-magnitude trials on each axis.
These sweeps were held out from inertia fitting, but their steady tails had
already been used for drag fitting.

Yaw response fitting gave $I_{\mathrm{eff}}=0.24254$ kg\,m$^2$, with
validation angular-rate RMSE 0.000928 rad/s. The dry yaw inertia is a geometry
estimate; the effective sum is the more directly validated quantity. These
are offline, scene-specific estimates, not physical-robot measurements or an
online adaptive identifier.

\section{Implementation limits}
The source does not enforce a wall-clock stop at the configured 80 ms
deadline, an explicit thruster-rate constraint, a terminal invariant set, or
robust recursive feasibility. Status 1 from L-BFGS-B is accepted when the
iterate and objective are finite; this is not a convergence certificate.
Historical timing and validation numbers above were read from saved reports,
not newly measured for this document.
\section{Equivalent condensed quadratic program}
This is an algebraic explanation, not a dense-matrix implementation.
Let $N=N_p$, $Z=[z_0^T,\ldots,z_{N-1}^T]^T$ and
$W=[w_0^T,\ldots,w_{N-1}^T]^T$. Repeated substitution gives
\begin{align}
z_i&=A_d^iz_0+\sum_{j=0}^{i-1}A_d^{i-1-j}(B_dw_j+\delta_j),\notag\\
Z&=\mathcal A z_0+\mathcal BW+\mathcal d=p+\mathcal BW,\notag\\
\mathcal A_i&=A_d^i,\qquad
\mathcal B_{ij}=\begin{cases}A_d^{i-1-j}B_d&j<i,\\0&j\ge i,\end{cases}\notag\\
\mathcal d_i&=\sum_{j=0}^{i-1}A_d^{i-1-j}\delta_j,
\qquad i,j\in\{0,\ldots,N-1\}.\tag{E26}
\end{align}
Empty sums are zero. Choosing E17 for $\delta$ recovers the current source;
choosing E16 yields the consistent correction-input model.
Let $L$ have identity diagonal blocks and negative identity first
subdiagonal blocks, so $LW$ stacks the increments with $\Delta w_0=w_0$.
Define
\[
\bar Q=\operatorname{blkdiag}(Q,\ldots,Q,F),\qquad
\bar R=I_N\otimes R,\qquad \bar S=I_N\otimes S.
\]
Then
\begin{align}
\mathcal J(W)&=(p+\mathcal BW)^T\bar Q(p+\mathcal BW)
 +W^T\bar RW+(LW)^T\bar S(LW)\notag\\
&=\tfrac12W^T\mathsf HW+g^TW+p^T\bar Qp,\notag\\
\mathsf H&=2(\mathcal B^T\bar Q\mathcal B+\bar R+L^T\bar SL),\qquad
g=2\mathcal B^T\bar Qp,\notag\\
\min_W\quad&\tfrac12W^T\mathsf HW+g^TW,\qquad \ell\le W\le u,\notag\\
\ell_i&=\underline a_i-a_{r,i},\quad u_i=\overline a_i-a_{r,i},\quad
\nabla_W\mathcal J=\mathsf HW+g.\tag{E27}
\end{align}
Since $R\succ0$ and the other quadratic terms are positive semidefinite,
$v^T\mathsf Hv\ge2v^T\bar Rv>0$ for every nonzero $v$.
Thus this frozen finite-dimensional problem is strictly convex. This does
not prove stability of the nonlinear closed loop. E20 computes the same
gradient without explicitly constructing the condensed matrices.
The last column block of $\mathcal B$ is zero: $w_{N-1}$ only affects the
unpenalized $z_N$. Its optimization still involves input and increment costs.

\section{One receding-horizon control cycle}
\begin{enumerate}
\item Convert measurements to ENU/FLU (E01), read the time-indexed reference,
and form $z_0$ (E14).
\item Shift the previous correction sequence (E21), predict with E17, and
recover nominal physical states using
$\hat\eta_i=\eta_{r,i}+\hat e_i$ and
$\hat\nu_i=J(\hat\psi_i)^T(\dot\eta_{r,i}+\widehat{\dot e}_i)$.
\item Compute E12 and solve E13 once per stage; freeze all 20 boxes.
\item Shift the boxes by reference acceleration (E18), then solve E19
with the joint objective and adjoint gradient E20. The current limit is
90 iterations; finite status 0/1 results are accepted.
\item Execute only $a_c=a_{r,0}+w_0^*$, apply E07, configured force scaling,
safety clipping, and the FLU-to-FRD force conversion.
\item The external allocator applies E23 and E24. Repeat using new
measurements at the next 0.1-second cycle.
\end{enumerate}
\end{document}
